\subsection{Data Preprocessing Steps}

As already mentioned, for the CICEVSE2024 dataset, it was only necessary to remove four rows containing \texttt{NaN} values to obtain clean and usable data for the initial processing step. \\

For the dataset generated from the experimental environment, the aim was first to include as much benign-status data as possible. Therefore, a PromQL request was made to retrieve monitoring data from a night when the environment was not used for other tests, as shown in Listing \ref{lst:prometheustocsv_scraper}. \\

Subsequently, the attacks were initiated to monitor the behavior of the HPC and kernel events in the environment under attack conditions, which can be seen in Listing \ref{lst:python_attack_script}.  
Table~\ref{tab:attacks} provides an overview of all attacks carried out during the experiment.

\begin{table}[H]
\centering
\begin{tabular}{|l|l|l|}
\hline
\textbf{Attack Type} & \textbf{Tool} & \textbf{Description / Parameters} \\
\hline
Slowloris          & Python script  & 500 sockets against target host \\
UDP-Flood          & hping3         & \texttt{--flood --udp -p 80} \\
ICMP-Flood         & hping3         & \texttt{--flood -1} \\
TCP-Flood          & hping3         & \texttt{--flood -p 80} \\
SYN-Flood          & hping3         & \texttt{--flood --rand-source -p 80} \\
Portscan           & nmap           & 1000 iterations, \texttt{-sT} scan \\
OS-Fingerprinting  & nmap           & 1000 iterations, \texttt{-O} scan \\
Aggressive-Scan    & nmap           & 1000 iterations, \texttt{-A} scan \\
\hline
\end{tabular}
\caption{Overview of executed attacks.}
\label{tab:attacks}
\end{table}

Using the generated log file, the correct attack intervals could be assigned to the dataset file for machine learning training and data analysis.  
Table~\ref{tab:intervals} shows the exact time periods used.

\begin{table}[H]
\centering
\begin{tabular}{|l|l|l|}
\hline
\textbf{Start Time} & \textbf{End Time} & \textbf{Attack Type} \\
\hline
2025-07-23 06:52:35 & 2025-07-23 07:07:35 & Slowloris \\
2025-07-23 07:27:35 & 2025-07-23 07:42:35 & UDP-Flood \\
2025-07-23 08:25:45 & 2025-07-23 08:47:50 & ICMP-Flood \\
2025-07-23 08:50:49 & 2025-07-23 09:05:49 & TCP-Flood \\
2025-07-23 09:05:49 & 2025-07-23 09:08:49 & SYN-Flood \\
2025-07-23 09:10:44 & 2025-07-23 09:20:44 & Portscan \\
2025-07-23 09:35:44 & 2025-07-23 09:45:44 & SYN-Flood \\
2025-07-23 10:00:44 & 2025-07-23 10:10:44 & OS-Fingerprinting \\
2025-07-23 10:25:44 & 2025-07-23 10:35:44 & Aggressive-Scan \\
2025-07-23 10:50:44 & 2025-07-23 11:00:44 & Slowloris \\
2025-07-23 11:10:21 & 2025-07-23 11:18:21 & UDP-Flood \\
2025-07-23 11:28:21 & 2025-07-23 11:36:21 & ICMP-Flood \\
2025-07-23 11:46:21 & 2025-07-23 11:54:21 & TCP-Flood \\
2025-07-23 12:08:55 & 2025-07-23 12:12:21 & Portscan \\
2025-07-23 12:22:21 & 2025-07-23 12:30:21 & SYN-Flood \\
\hline
\end{tabular}
\caption{Attack intervals during data collection.}
\label{tab:intervals}
\end{table}

These preprocessing steps enabled a meaningful exploratory data analysis, which allowed for the identification of key features for training machine learning models to achieve optimal results.  
This step is essential, as most models do not perform well with 900 or more features. This is also demonstrated in source \cite{EV_Security}, where the best results were achieved using only 11 to 20 key features. \\

Moreover, using too many features would lead to performance issues on the Raspberry Pi in the live system, as it would be required to query and process values for up to 900 features every 5 to 60 seconds.

% Insert figure demonstrating performance issues when querying 900 features
